\section{Spin and Fermions}

The sites carry spin at the base. Matter is fermionic. The knit keeps matter stable against collapse. None of
the three facts is added to the model. Each one is read off the eight atoms already in hand, the way the $24$
sites of a dock transform under their own symmetry, the topology of a twist in the tone field, and the measured
sign of one piece of the knit's collide step. This section reads them off in order.

Two terms recur throughout, so a plain statement of each helps. A \emph{spinor} is the mathematical object that
describes a piece of matter with half-integer spin. Its strange signature is that one full turn does not bring it
home. Rotate a spinor through $360$ degrees and it picks up a minus sign. Only a second full turn, $720$ degrees
in all, restores it. This is the \emph{double cover}, and it is the defining mark of half-integer spin. A
\emph{fermion} is a piece of matter, an electron or a quark, no two of which can occupy the same state. That
exclusion is why matter takes up space at all. Fermions have half-integer spin, so every fermion is described by
a spinor. The two ideas are bound together, and the substrate delivers that binding rather than presupposing it.

\begin{result}[The matter sector is read off the sites]
The $24$ sites of a dock are not just $24$ arrows. Their split under their own symmetry is a vector plus two
spinors, locked together. That split is the seed of all matter.
\end{result}

\subsection{The triality split of the sites}

Recall the base picture. A dock is a here-now, and out of each dock run exactly $24$ sites, the $24$ directions
of the $\{3,4,3,4\}$ honeycomb. There is no rest slot, no twenty-fifth site. Those $24$ directions are the
vertices of the $24$-cell, equivalently the $D_4$ root system. The whole spin story turns on how that set of $24$
transforms under its own symmetry group.

Under $\mathrm{SO}(8)$ the $24$ sites decompose as three sets of eight,
\[
24 = 8v + 8s + 8c.
\]
The summand $8v$ is the $\mathrm{SO}(8)$ vector, of integer spin. The two summands $8s$ and $8c$ are genuine
half-integer-spin representations, the two spinors of $\mathrm{SO}(8)$. The outer symmetry that permutes the
three eights is \emph{triality}, the order-three $S_3$ symmetry of the $D_4$ Dynkin diagram. Triality is the
exceptional feature of $D_4$ alone, and it is the reason the vector and the two spinors sit on an equal footing.

\begin{table}[ht]
\centering
\begin{tabular}{@{}llll@{}}
\toprule
sector & name & spin & role \\
\midrule
$8v$ & the vector & integer & the photon \\
$8s$ & a spinor & half-integer & one fermion chirality \\
$8c$ & the other spinor & half-integer & the other fermion chirality \\
\bottomrule
\end{tabular}
\caption{The triality split of the $24$ sites under $\mathrm{SO}(8)$, with triality $S_3$ permuting the three
eights.}
\label{tab:triality-split}
\end{table}

\begin{proposition}[The spinor decomposition]
The $24$ sites of a dock, as the $D_4$ root system, decompose under $\mathrm{SO}(8)$ as $8v + 8s + 8c$. The
summands $8s$ and $8c$ are half-integer-spin representations, so the substrate carries spinors directly. Triality
permutes the three eights and is the seed of the family structure.
\end{proposition}

The force of this is that spin is not glued on. It is the shape of the sites themselves. Two of the three pieces
of a dock are already spinors before any dynamics runs. The substrate carries spin \textbf{at the base}.

This split is exactly what the structural co-emergence result reads off the rotation subgroup. Every rotation of
the dock keeps each of $8v$, $8s$, and $8c$ as an invariant subspace, while triality relates the three to one
another. That is a structural fact about the group, true before any beat ticks. It is the static backbone on
which the dynamical coupling of the photon and fermion sectors, taken up later, is hung.

\subsection{The control: non-spinor sites}

The claim that the sites carry spin is only sharp if some plausible alternative fails to. The discriminator is
what happens on a substrate whose sites are not $D_4$.

Take the $12$ directions of the $\{5,3,4\}$ honeycomb, the curved substrate of an earlier chapter. Under its
icosahedral symmetry the $12$ directions decompose as
\[
12 = 1 + 3 + 3' + 5,
\]
a singlet, two distinct triplets, and a five. Every piece on that list is integer spin. There is no spinor
anywhere in it. A substrate built on those $12$ directions would carry no fundamental spinor at all.

\begin{result}[The spinor gap selects the substrate]
The $12$ sites of $\{5,3,4\}$ carry no spinor. The $24$ sites of $\{3,4,3,4\}$ carry two. That gap is why
$\{3,4,3,4\}$ is the matter substrate.
\end{result}

This is the load-bearing selection argument for the matter sector. To host fundamental fermion fields a
substrate must carry the $8s$ and $8c$ spinors among its sites. Only the $24$ $D_4$ sites do. The full forcing,
set beside the crystallographic requirements, is developed where the tessellation survey and the dimension
ladder lay the candidate honeycombs side by side.

\subsection{A self is automatically a fermion}

There is a second route to spin, independent of the spinor sites, and it does not need them at all. It rests only
on topology.

A \emph{self}, in the sense used through the paper, is a persistent self-maintaining knot in the tone field, a
twist woven across many docks over many beats that cannot be combed out. The matter object the theory identifies
with a fundamental fermion is one particular such knot, a \emph{hopfion}.

\begin{definition}[Matter soliton]
A \emph{matter soliton} is a topological knot in the tone field, a \emph{hopfion}, classified by its Hopf charge
in $\pi_3(S^2) = \mathbb{Z}$. The integer charge counts how the tone configuration links with itself. It cannot
change under any smooth evolution of the field.
\end{definition}

A hopfion of odd Hopf charge picks up a minus sign under a $2\pi$ rotation. This is the Finkelstein-Rubinstein
result, sharpened by the spin-statistics argument of Wilczek and Zee \cite{finkelsteinrubinstein1968,
wilczekzee1983}. A minus sign under a single full turn is the defining mark of a fermion. So the spin-statistics
connection here is geometric. A self of odd charge is automatically a fermion, with nothing assumed about its
statistics beyond its topology. In the carrying image such a self is a braid in the fabric, a knot of fire, and
the weave itself decides the statistics.

\begin{proposition}[Odd-charge selves are fermions]
A topological soliton in the tone field with odd Hopf charge transforms with a sign change under a $2\pi$
rotation. By the spin-statistics connection it is a fermion. The matter solitons are therefore fermions as a
consequence of their topology, with no separate statistical postulate.
\end{proposition}

\subsection{Two sources of spin, kept distinct}

Spin enters this theory in two ways, and they answer different questions. Keeping them apart is essential,
because the selection claim depends on the distinction being exact.

The first source is the \emph{spinor sites}, the $D_4$ pieces $8s$ and $8c$. These are the fundamental fermion
fields. They genuinely require $\{3,4,3,4\}$, since the $12$-direction alternative carries no spinor at all. The
second source is the \emph{topological route}, the identification of a self with a hopfion and hence with a
fermion. That route is substrate-independent. It works even on $\{5,3,4\}$, because it rests only on
$\pi_3(S^2)$ and not on the direction set.

\begin{table}[ht]
\centering
\begin{tabular}{@{}lll@{}}
\toprule
source & what it is & substrate need \\
\midrule
the spinor sites & the $D_4$ $8s + 8c$, fundamental fermion fields & genuinely needs $\{3,4,3,4\}$ \\
the topological route & self $=$ hopfion $=$ fermion & works on any substrate, even $\{5,3,4\}$ \\
\bottomrule
\end{tabular}
\caption{The two ways spin appears, and what each one requires of the substrate.}
\label{tab:two-spin-sources}
\end{table}

The two agree, but the precise selection claim is narrow and must be stated exactly.

\begin{principle}[What forces the substrate]
The \emph{spinor sites} force $\{3,4,3,4\}$. It is not that spin of any kind forces it. The fundamental matter
fields and the $\mathrm{SO}(10)$ gauge story need the $8s$ and $8c$ that only the $24$ $D_4$ sites carry. The
topological route is the separate reason emergent selves come out fermionic on their own.
\end{principle}

\subsection{Why matter forms and persists}

A spinor among the sites is only potential. Whether a self actually holds together rather than collapsing or
spreading is a separate question, settled by Derrick's theorem.

Rescale a soliton by a factor $\lambda$. In $d$ spatial dimensions the exchange energy scales as $\lambda^{d-2}$,
while a higher-gradient Skyrme term scales as $\lambda^{d-4}$. In three dimensions exchange alone scales as
$\lambda$, so the soliton shrinks without bound and collapses to a point. A higher-gradient Skyrme term, scaling
as $\lambda^{-1}$, is required to arrest the collapse \cite{skyrme1962}. In the two-dimensional cusp of the
curved substrate exchange is scale-marginal and no stabilizer is needed, a clean tie to the bulk-cusp dimension
split. In the three-dimensional flat cusp of $\{3,4,3,4\}$ a stabilizer is mandatory.

The Derrick squeeze is confirmed directly. On a real three-dimensional texture the exchange energy grows with
size while the Skyrme energy falls, exactly the competition Derrick predicts. So the question of stability reduces
to a single sign. Does the knit's collide step supply a Skyrme term of the stabilizing sign? The collide step is
the knit's last free part, fixed only by being conserving, reversible, local, and lattice-symmetric. The Skyrme
coefficient appears at the higher-gradient, Burnett order of the collide step's gradient expansion. Stability is
therefore a definite sign of a definite quantity.

The honest record on this sign is mixed, and is stated plainly rather than papered over.

\begin{itemize}
\item The soliton-binding experiment shows two solitons binding at a finite separation, with the rest mass adding
up with topological charge. Matter behaves like matter.
\item The quaternion-twist experiment shows that a pure quaternion twist, the handedness of the sites, binds a
stable Skyrmion with no extra coupling needed.
\item The actual-knit experiment shows that a one-dimensional fermion sea does not on its own settle the
three-dimensional stabilizing sign. This is logged as an honest open gate.
\end{itemize}

\begin{result}[Where the sign is settled, matter forms]
Where the three-dimensional stabilizing sign is settled, matter forms. Where it is not yet settled, the record
says so. Even in the worst case, matter-formation simply selects the collide member with the stabilizing sign.
The base needs no fine-tuning for matter to exist.
\end{result}

\subsection{What is base and what is emergent}

The base is discrete and deterministic. The line between what is given and what emerges is sharp here, and worth
drawing.

The $24$ sites and their $D_4$ spinor split are base. The minus-sign double cover is a fact about the group, also
base. A self is an emergent twist in the tone field, a mid-layer object woven across many docks over many beats.
Its fermionic statistics emerge from its topology, not from any added piece of the knit. The complex quantum
amplitudes that name spin states are the emergent description of the real discrete walk. The base is the walk.

\begin{table}[ht]
\centering
\begin{tabular}{@{}ll@{}}
\toprule
ingredient & status \\
\midrule
the $24$ sites and the $8v + 8s + 8c$ split & base \\
the $2\pi$ double-cover minus sign & base, a fact about the group \\
a self, a knot in the tone field & emergent, a mid-layer object \\
the fermionic statistics of a self & emergent, from topology \\
the complex spin amplitudes & emergent description of the discrete walk \\
\bottomrule
\end{tabular}
\caption{The base ingredients of spin against the emergent ones.}
\label{tab:base-vs-emergent}
\end{table}

The matter sector is not bolted on. The sites carry two spinors, so matter is spinorial. A self is automatically
a fermion, so statistics are geometric. The knit supplies the stabilizer where it has been measured. The
half-integer spin that defines matter is already present in the way a dock's $24$ sites turn under their own
symmetry.

\begin{result}[Fermions are the stable knots]
Fermions are not added to the model. They are the stable knots of the fire, and the fire is the sites turning
twice.
\end{result}
